% GENERATED FILE. Edit canonical lessons, then run npm run generate:books.
% canonical-source-hash: fnv1a64:61d9e924fffb7ca3
% canonical-lessons: FR-C424-auberge, FR-C424-gite, FR-C424-nuitee, FR-C424-loger, FR-C424-sejourner

\chapter{Où dormir ?}
\label{ch:fr-ou-dormir}
\hlchaptermodality{424}

\begin{chapteropening}
\textbf{By the end of this chapter:} \emph{I can talk about places to stay: the inn, the holiday cottage, a night, staying somewhere and spending time there.}

Complete the last lesson of chapter 424: I can talk about places to stay: the inn, the holiday cottage, a night, staying somewhere and spending time there.
\end{chapteropening}

\section[l'auberge]{l'auberge — the inn; the hostel}

\label{lesson:FR-C424-auberge}



\begin{quote}
Before the new one: say the French for to capsize, then the French for to drift.
\end{quote}

\subsection*{You'll want to know: l'auberge}
\textbf{l'auberge} — \textquotedblleft{}the inn; the hostel\textquotedblright{}.

\begin{quote}
Nous avons dormi dans une petite auberge de campagne.
\end{quote}

\begin{quote}
\emph{We slept in a little country inn.}
\end{quote}

A feminine noun. \emph{Une auberge de jeunesse} — a youth hostel.

\begin{cousinweb}
\emph{auberge} was borrowed from Provençal \emph{alberga}, which is usually traced to a Germanic word meaning 'shelter for an army'. EtymDB lists it without a source.
\end{cousinweb}

\begin{grammarlens}[title={What you've built}]
One more word for telling, explaining and describing at B1.
\end{grammarlens}

\subsection*{Your turn}
\begin{itemize}
\raggedright
  \item \emph{Say it:} \emph{l'auberge}
  \item \emph{Say it:} \emph{l'auberge}, once more
  \item \emph{Say it:} a sentence of your own with \emph{l'auberge}
  \item \emph{Recall:} say the French for to capsize, then the French for to drift, then say \emph{l'auberge} again
\end{itemize}

\begin{tcolorbox}[breakable,colback=teal!4,colframe=teal!35!black,title={Before you move on}]
What does l'auberge mean? (\textquotedblleft{}The inn; the hostel\textquotedblright{}.) Say it once more.
\end{tcolorbox}
\section[le gîte]{le gîte — the holiday cottage; the shelter}

\label{lesson:FR-C424-gite}



\begin{quote}
Before the new one: say the French for to drift, then the French for the inn.
\end{quote}

\subsection*{You'll want to know: le gîte}
\textbf{le gîte} — \textquotedblleft{}the holiday cottage; the shelter\textquotedblright{}.

\begin{quote}
Cet été, ils louent un gîte à la campagne.
\end{quote}

\begin{quote}
\emph{This summer they are renting a holiday cottage in the country.}
\end{quote}

\emph{Un gîte rural} — a holiday cottage in the country.

\begin{cousinweb}
\emph{gîte} comes from Old French \emph{giste} (EtymDB), built on \emph{gésir} 'to lie down', which is inherited from Latin \emph{iacēre} 'to lie' (Lewis \& Short).
\end{cousinweb}

\begin{grammarlens}[title={What you've built}]
One more word for telling, explaining and describing at B1.
\end{grammarlens}

\subsection*{Your turn}
\begin{itemize}
\raggedright
  \item \emph{Say it:} \emph{le gîte}
  \item \emph{Say it:} \emph{le gîte}, once more
  \item \emph{Say it:} a sentence of your own with \emph{le gîte}
  \item \emph{Recall:} say the French for to drift, then the French for the inn, then say \emph{le gîte} again
\end{itemize}

\begin{tcolorbox}[breakable,colback=teal!4,colframe=teal!35!black,title={Before you move on}]
What does le gîte mean? (\textquotedblleft{}The holiday cottage; the shelter\textquotedblright{}.) Say it once more.
\end{tcolorbox}
\section[la nuitée]{la nuitée — the overnight stay (paid for at a hotel)}

\label{lesson:FR-C424-nuitee}



\begin{quote}
Before the new one: say the French for the inn, then the French for the holiday cottage.
\end{quote}

\subsection*{You'll want to know: la nuitée}
\textbf{la nuitée} — \textquotedblleft{}the overnight stay (paid for at a hotel)\textquotedblright{}.

\begin{quote}
La nuitée coûte quatre-vingts euros avec le petit-déjeuner.
\end{quote}

\begin{quote}
\emph{A night costs eighty euros with breakfast.}
\end{quote}

A word from hotel bills and prices.

\begin{cousinweb}
\emph{nuitée} is built on \emph{nuit} 'night' (EtymDB), which is inherited from Latin \emph{nox, noctis} 'night'.
\end{cousinweb}

\begin{grammarlens}[title={What you've built}]
One more word for telling, explaining and describing at B1.
\end{grammarlens}

\subsection*{Your turn}
\begin{itemize}
\raggedright
  \item \emph{Say it:} \emph{la nuitée}
  \item \emph{Say it:} \emph{la nuitée}, once more
  \item \emph{Say it:} a sentence of your own with \emph{la nuitée}
  \item \emph{Recall:} say the French for the inn, then the French for the holiday cottage, then say \emph{la nuitée} again
\end{itemize}

\begin{tcolorbox}[breakable,colback=teal!4,colframe=teal!35!black,title={Before you move on}]
What does la nuitée mean? (\textquotedblleft{}The overnight stay (paid for at a hotel)\textquotedblright{}.) Say it once more.
\end{tcolorbox}
\section[loger]{loger — to lodge, to stay; to put up}

\label{lesson:FR-C424-loger}



\begin{quote}
Before the new one: say the French for the holiday cottage, then the French for the overnight stay.
\end{quote}

\subsection*{You'll want to know: loger}
\textbf{loger} — \textquotedblleft{}to lodge, to stay; to put up\textquotedblright{}.

\begin{quote}
Pendant le festival, nous logeons chez des amis.
\end{quote}

\begin{quote}
\emph{During the festival we are staying with friends.}
\end{quote}

\emph{Le logement} — housing.

\begin{cousinweb}
\emph{loger} comes from Old French \emph{logier} (EtymDB), built on \emph{loge} 'hut', which is usually traced to a Frankish word for a shelter of leaves.
\end{cousinweb}

\begin{grammarlens}[title={What you've built}]
One more word for telling, explaining and describing at B1.
\end{grammarlens}

\subsection*{Your turn}
\begin{itemize}
\raggedright
  \item \emph{Say it:} \emph{loger}
  \item \emph{Say it:} \emph{loger}, once more
  \item \emph{Say it:} a sentence of your own with \emph{loger}
  \item \emph{Recall:} say the French for the holiday cottage, then the French for the overnight stay, then say \emph{loger} again
\end{itemize}

\begin{tcolorbox}[breakable,colback=teal!4,colframe=teal!35!black,title={Before you move on}]
What does loger mean? (\textquotedblleft{}To lodge, to stay; to put up\textquotedblright{}.) Say it once more.
\end{tcolorbox}
\section[séjourner]{séjourner — to stay for a while, to spend time (somewhere)}

\label{lesson:FR-C424-sejourner}



\begin{quote}
Before the new one: say the French for the overnight stay, then the French for to lodge.
\end{quote}

\subsection*{You'll want to know: séjourner}
\textbf{séjourner} — \textquotedblleft{}to stay for a while, to spend time (somewhere)\textquotedblright{}.

\begin{quote}
Ils ont séjourné une semaine au bord du lac.
\end{quote}

\begin{quote}
\emph{They stayed a week by the lake.}
\end{quote}

More formal than \emph{rester}.

\begin{cousinweb}
\emph{séjourner} goes back to Vulgar Latin \emph{subdiurnāre} 'to spend the day', built on Latin \emph{diurnus} 'of the day', from \emph{diēs} 'day' (Lewis \& Short).
\end{cousinweb}

\begin{grammarlens}[title={What you've built}]
One more word for telling, explaining and describing at B1.
\end{grammarlens}

\subsection*{Your turn}
\begin{itemize}
\raggedright
  \item \emph{Say it:} \emph{séjourner}
  \item \emph{Say it:} \emph{séjourner}, once more
  \item \emph{Say it:} a sentence of your own with \emph{séjourner}
  \item \emph{Recall:} say the French for the overnight stay, then the French for to lodge, then say \emph{séjourner} again
\end{itemize}

\begin{tcolorbox}[breakable,colback=teal!4,colframe=teal!35!black,title={Before you move on}]
What does séjourner mean? (\textquotedblleft{}To stay for a while, to spend time (somewhere)\textquotedblright{}.) Say it once more.
\end{tcolorbox}
